\documentclass[10pt,a4paper]{amsart}
\usepackage{amsmath,amssymb,amsthm,graphicx}
% Retained preamble for the Stacks Project chapter morphisms.tex (Morphisms of Schemes).
% The chapter file pinned-original-source.tex carries only \input{preamble} on its first line, so no
% preamble text lives inside the chapter itself.  The Stacks Project uses one single shared preamble for
% every chapter; the snapshot stages it verbatim in sources/stacks-chow/preamble.tex.  The block below is
% copied mechanically, line for line, from that file's lines 14-18 (the Xy-pic setup) and lines 154-217
% (the native theorem declarations together with the mathematical macros).  Nothing was added, removed or
% redefined.  The wrapper class and the amsmath/amssymb/amsthm/graphicx base packages are supplied by the
% assembler, so the class/\externaldocument part of the shared preamble is not repeated here.
% For commutative diagrams we use Xy-pic
\usepackage[all]{xy}

% We use 2cell for 2-commutative diagrams.
\xyoption{2cell}
% Theorem environments.
%
\theoremstyle{plain}
\newtheorem{theorem}[subsection]{Theorem}
\newtheorem{proposition}[subsection]{Proposition}
\newtheorem{lemma}[subsection]{Lemma}

\theoremstyle{definition}
\newtheorem{definition}[subsection]{Definition}
\newtheorem{example}[subsection]{Example}
\newtheorem{exercise}[subsection]{Exercise}
\newtheorem{situation}[subsection]{Situation}

\theoremstyle{remark}
\newtheorem{remark}[subsection]{Remark}
\newtheorem{remarks}[subsection]{Remarks}

\numberwithin{equation}{subsection}

% Macros
%
\def\lim{\mathop{\mathrm{lim}}\nolimits}
\def\colim{\mathop{\mathrm{colim}}\nolimits}
\def\Spec{\mathop{\mathrm{Spec}}}
\def\Hom{\mathop{\mathrm{Hom}}\nolimits}
\def\Ext{\mathop{\mathrm{Ext}}\nolimits}
\def\SheafHom{\mathop{\mathcal{H}\!\mathit{om}}\nolimits}
\def\SheafExt{\mathop{\mathcal{E}\!\mathit{xt}}\nolimits}
\def\Sch{\mathit{Sch}}
\def\Mor{\mathop{\mathrm{Mor}}\nolimits}
\def\Ob{\mathop{\mathrm{Ob}}\nolimits}
\def\Sh{\mathop{\mathit{Sh}}\nolimits}
\def\NL{\mathop{N\!L}\nolimits}
\def\CH{\mathop{\mathrm{CH}}\nolimits}
\def\proetale{{pro\text{-}\acute{e}tale}}
\def\etale{{\acute{e}tale}}
\def\QCoh{\mathit{QCoh}}
\def\Ker{\mathop{\mathrm{Ker}}}
\def\Im{\mathop{\mathrm{Im}}}
\def\Coker{\mathop{\mathrm{Coker}}}
\def\Coim{\mathop{\mathrm{Coim}}}

% Boxtimes
%
\DeclareMathSymbol{\boxtimes}{\mathbin}{AMSa}{"02}

%
% Macros for moduli stacks/spaces
%
\def\QCohstack{\mathcal{QC}\!\mathit{oh}}
\def\Cohstack{\mathcal{C}\!\mathit{oh}}
\def\Spacesstack{\mathcal{S}\!\mathit{paces}}
\def\Quotfunctor{\mathrm{Quot}}
\def\Hilbfunctor{\mathrm{Hilb}}
\def\Curvesstack{\mathcal{C}\!\mathit{urves}}
\def\Polarizedstack{\mathcal{P}\!\mathit{olarized}}
\def\Complexesstack{\mathcal{C}\!\mathit{omplexes}}
% \Pic is the operator that assigns to X its picard group, usage \Pic(X)
% \Picardstack_{X/B} denotes the Picard stack of X over B
% \Picardfunctor_{X/B} denotes the Picard functor of X over B
\def\Pic{\mathop{\mathrm{Pic}}\nolimits}
\def\Picardstack{\mathcal{P}\!\mathit{ic}}
\def\Picardfunctor{\mathrm{Pic}}
\def\Deformationcategory{\mathcal{D}\!\mathit{ef}}

% --- appended 2026-10-12 (Chief): the project own definitions of the reference/slogan/history
% display environments, copied mechanically from the shared Stacks preamble lines 8-12
% (\usepackage{verbatim} plus the three \newenvironment lines, which the chapter preamble keeps
% verbatim from the project preamble). Nothing added, removed or redefined: these environments
% carry no mathematical content and the project itself typesets them as comments. Without these
% three lines a counted statement whose own text uses them fails with Environment slogan undefined.
\usepackage{verbatim}
\newenvironment{reference}{\comment}{\endcomment}
\newenvironment{slogan}{\comment}{\endcomment}
\newenvironment{history}{\comment}{\endcomment}
\usepackage{cleveref}
\crefname{originallemma}{original lemma}{original lemmas}
\Crefname{originallemma}{Original lemma}{Original lemmas}
\makeatletter
\newlabel{lemma-factor-quasi-compact-immersion}{{01QV}{source}{}{sourcepointer-1}{}}
\newlabel{lemma-factor-quasi-compact-immersion@cref}{{[originallemma][0][]01QV}{[1][0][]source}{}{}{}}
\makeatother
\crefname{originallemma}{original lemma}{original lemmas}
\Crefname{originallemma}{Original lemma}{Original lemmas}
\makeatletter
\newlabel{lemma-factor-reduced-immersion}{{03DQ}{source}{}{sourcepointer-2}{}}
\newlabel{lemma-factor-reduced-immersion@cref}{{[originallemma][0][]03DQ}{[1][0][]source}{}{}{}}
\makeatother
\crefname{originallemma}{original lemma}{original lemmas}
\Crefname{originallemma}{Original lemma}{Original lemmas}
\makeatletter
\newlabel{lemma-quasi-compact-scheme-theoretic-image}{{01R8}{source}{}{sourcepointer-3}{}}
\newlabel{lemma-quasi-compact-scheme-theoretic-image@cref}{{[originallemma][0][]01R8}{[1][0][]source}{}{}{}}
\makeatother
\title{An immersion that is quasi-compact or has reduced source is an open immersion onto a scheme theoretically dense open of its scheme theoretic image}
\author{The Stacks Project Authors}
\date{stacks-project:morphisms.tex:local-snapshot-20261012; GNU FDL 1.2 or later}
\begin{document}
\maketitle
\noindent\textit{Editorial note.} All delivered mathematical text is the author's own, copied from the pinned Stacks Project chapter file pinned-original-source.tex; no statement or proof step was rewritten, reordered or removed.\par\medskip
\noindent\textit{Editorial note.} Following the multi-claim convention for this source, source\_url is the source-level URL https://stacks.math.columbia.edu/ and the source label lemma-quasi-compact-immersion is carried in claim\_key and here rather than in a per-claim tag URL.\par\medskip
\noindent\textit{Editorial note.} The preamble is the chapter's shared Stacks Project preamble, retained as sources/stacks-morphisms/retained-preamble.tex: a line-for-line excerpt of the shared Stacks preamble staged with this snapshot. No macro was invented or redefined, and this card adds no preamble text of its own.\par\medskip
\noindent\textit{Editorial note.} Statement span 1138-1146 and proof span 1148-1170 come from the halfA pool card operations/pool-cards/stacks-morphisms-halfA.json and were re-read in the pinned chapter file before this card was written.\par\medskip
\noindent\textit{Editorial note.} The proof distinguishes the two hypotheses of the statement with the author's phrases 'In the case that Z is reduced' and 'In the case that Z to X is quasi-compact'; both are preserved.\par\medskip
\noindent\textit{Editorial note.} 3 cross-references occur inside the counted spans; they are listed in external\_references as navigation pointers only.\par\medskip
\noindent\textit{Allowed standard prerequisites (editorial).}\begin{itemize}
\item The factorization of an immersion as an open immersion followed by a closed immersion, in the quasi-compact case and in the reduced case, cited from the same chapter. Cited at proof lines 3365-3367. Produces the closed intermediate scheme Z-bar-one inside which the source sits as an open subscheme.
\item The scheme theoretic image of a morphism as the smallest closed subscheme through which it factors, and its density properties in the quasi-compact case, cited from the same chapter. Cited at proof lines 3368-3370 and 3384-3387. Places the scheme theoretic image inside the intermediate closed subscheme and supplies scheme theoretic and topological density.
\item A closed immersion into a reduced scheme that is a bijection on underlying topological spaces is an isomorphism. Used at proof lines 3377-3381. Identifies the intermediate closed subscheme with the scheme theoretic image in the reduced case.
\end{itemize}\medskip
\ifdefined\hypertarget\hypertarget{sourcepointer-1}{}\fi
\noindent\textit{Source reference.} Original lemma 01QV: focused-seeds/AGP/sources/stacks-snapshot-20261011/tags.txt, source lines 2280--2280. Same-chapter reference used in the retained author text at source line(s) 1149. It is supplied only as a navigation pointer so that the retained text compiles; the counted item delivers the author's own statement and the author's complete proof, and no part of the referenced item is copied in or substituted for the counted argument.\par\medskip
\ifdefined\hypertarget\hypertarget{sourcepointer-2}{}\fi
\noindent\textit{Source reference.} Original lemma 03DQ: focused-seeds/AGP/sources/stacks-snapshot-20261011/tags.txt, source lines 4429--4429. Same-chapter reference used in the retained author text at source line(s) 1150, 1159. It is supplied only as a navigation pointer so that the retained text compiles; the counted item delivers the author's own statement and the author's complete proof, and no part of the referenced item is copied in or substituted for the counted argument.\par\medskip
\ifdefined\hypertarget\hypertarget{sourcepointer-3}{}\fi
\noindent\textit{Source reference.} Original lemma 01R8: focused-seeds/AGP/sources/stacks-snapshot-20261011/tags.txt, source lines 2293--2293. Same-chapter reference used in the retained author text at source line(s) 891, 903, 920, 928, 1007, 1084, 1167, 1169, 1369, 4982, 10102, 11100, 14144. It is supplied only as a navigation pointer so that the retained text compiles; the counted item delivers the author's own statement and the author's complete proof, and no part of the referenced item is copied in or substituted for the counted argument.\par\medskip
\begin{lemma}
\label{lemma-quasi-compact-immersion}
Let $h : Z \to X$ be an immersion. Assume either $h$ is quasi-compact
or $Z$ is reduced. Let $\overline{Z} \subset X$ be the scheme theoretic
image of $h$. Then the morphism $Z \to \overline{Z}$ is an open immersion
which identifies $Z$ with a scheme theoretically dense open
subscheme of $\overline{Z}$. Moreover, $Z$ is topologically
dense in $\overline{Z}$.
\end{lemma}

\begin{proof}
By Lemma \ref{lemma-factor-quasi-compact-immersion} or
Lemma \ref{lemma-factor-reduced-immersion} we can factor
$Z \to X$ as $Z \to \overline{Z}_1 \to X$ with $Z \to \overline{Z}_1$
open and $\overline{Z}_1 \to X$ closed. On the other hand, let
$Z \to \overline{Z} \subset X$ be the scheme theoretic image of
$Z \to X$. We conclude that $\overline{Z} \subset \overline{Z}_1$.
Since $Z$ is an open subscheme of $\overline{Z}_1$ it follows
that $Z$ is an open subscheme of $\overline{Z}$ as well.
In the case that $Z$ is reduced we know that $Z \subset \overline{Z}_1$
is topologically dense by the construction of $\overline{Z}_1$ in
the proof of Lemma \ref{lemma-factor-reduced-immersion}.
Hence $\overline{Z}_1$ and $\overline{Z}$ have the same
underlying topological spaces. Thus $\overline{Z} \subset \overline{Z}_1$
is a closed immersion into a reduced scheme which induces a bijection
on underlying topological spaces, and hence it is an isomorphism.
In the case that $Z \to X$ is quasi-compact we argue as follows:
The assertion that $Z$ is scheme theoretically dense in
$\overline{Z}$ follows from
Lemma \ref{lemma-quasi-compact-scheme-theoretic-image} part (3).
The last assertion follows from
Lemma \ref{lemma-quasi-compact-scheme-theoretic-image} part (4).
\end{proof}

\end{document}
